\documentclass[11pt]{article}

% Aktuelles offizielles *ACL-Layout. acl.sty und acl_natbib.bst unverändert
% im selben Projektverzeichnis belassen. Für diese nicht-anonyme Seminarabgabe
% wird der offizielle Final-Modus verwendet.
\usepackage{acl}
\usepackage{times}
\usepackage{latexsym}
\usepackage[T1]{fontenc}
\usepackage[utf8]{inputenc}
\usepackage[ngerman]{babel}
\usepackage{microtype}
\usepackage{inconsolata}
\usepackage{graphicx}
\usepackage{amsmath,amssymb}
\usepackage{booktabs}
\usepackage{tcolorbox}
\tcbuselibrary{skins,breakable}

% Deutsche Bezeichnungen auch dort erzwingen, wo acl.sty englische Überschriften fest vorgibt.
\usepackage{etoolbox}
\makeatletter
\AtBeginDocument{%
  \renewcommand{\abstractname}{Zusammenfassung}%
  \patchcmd{\thebibliography}
    {\section*{References\@mkboth {References}{References}}}
    {\section*{Literaturverzeichnis\@mkboth {Literaturverzeichnis}{Literaturverzeichnis}}}
    {}{}%
}
\makeatother

% BLUEPRINT-HINWEIS: Die folgenden tcolorbox-Umgebungen dienen nur als Arbeitshilfen.
% Sie können während des Schreibens bestehen bleiben, sollen aber vor der finalen
% Seminarabgabe in Fließtext, Tabellen oder Kommentare überführt bzw. entfernt werden.
% ACL lädt geometry, xcolor, natbib, hyperref und den Stil acl_natbib bereits selbst.
% Diese Pakete bzw. Stile hier nicht zusätzlich laden oder überschreiben.

% Farbdefinitionen und Box-Stile für das wissenschaftliche Arbeitsgerüst
\definecolor{boxblue}{RGB}{235, 243, 250}
\definecolor{borderblue}{RGB}{41, 128, 185}
\definecolor{boxgreen}{RGB}{234, 250, 234}
\definecolor{bordergreen}{RGB}{39, 174, 96}
\definecolor{boxorange}{RGB}{254, 249, 231}
\definecolor{borderorange}{RGB}{211, 84, 0}

\newtcolorbox{leitfragenbox}[1][]{
    enhanced,
    breakable,
    colback=boxblue,
    colframe=borderblue,
    fonttitle=\bfseries,
    title={Leitfragen f\"ur diesen Abschnitt},
    arc=2mm,
    boxrule=0.8pt,
    left=2mm, right=2mm, top=1.5mm, bottom=1.5mm,
    #1
}

\newtcolorbox{quellenbox}[1][]{
    enhanced,
    breakable,
    colback=boxgreen,
    colframe=bordergreen,
    fonttitle=\bfseries,
    title={Empfohlene Quellen \& Zitationszweck},
    arc=2mm,
    boxrule=0.8pt,
    left=2mm, right=2mm, top=1.5mm, bottom=1.5mm,
    #1
}

\newtcolorbox{faktenbox}[1][]{
    enhanced,
    breakable,
    colback=boxorange,
    colframe=borderorange,
    fonttitle=\bfseries,
    title={Verbindliche experimentelle Eckdaten / offene Verifikationen},
    arc=2mm,
    boxrule=0.8pt,
    left=2mm, right=2mm, top=1.5mm, bottom=1.5mm,
    #1
}

\title{Vergleich von Textrepräsentationen und Chunk-Selektion mit Enthaltung beim Belegstellen-Retrieval in deutschen Prüfungsordnungen}

% TODO vor der Abgabe: Platzhalter durch die endgültigen Autorendaten ersetzen.
\author{Seminarteilnehmer / Verfasser \\
  Institut f\"ur Sprache und Information \\
  Heinrich-Heine-Universit\"at D\"usseldorf \\
  \texttt{TODO: Universitäts-E-Mail-Adresse}}

\begin{document}

\maketitle

% BLUEPRINT: Zusammenfassung erst nach dem Einfrieren des finalen Fließtexts endgültig ausarbeiten.
\begin{abstract}
Diese Arbeit untersucht das Retrieval von Belegstellen (Chunks) in einer hierarchisch strukturierten deutschen Prüfungsordnung und berücksichtigt dabei eine explizite Enthaltungsentscheidung für Fragen, die sich aus dem Quelldokument nicht beantworten lassen. Der projektspezifische Datensatz umfasst 720 Fragen, darunter 120 absichtlich unbeantwortbare Zero-Gold-Fragen, sowie 201 Retrieval-Chunks. Verglichen werden eine sparse TF-IDF-Repräsentation und zwei feste multilinguale dichte Repräsentationen, Sentence-BERT und multilingual E5. Zusätzlich werden vier Familien von Chunk-Selektionsverfahren evaluiert: Top-$k$, absoluter Schwellenwert, Top-$k$ mit absolutem Schwellenwert und Auswahl über eine relative Wertemarge. Hyperparameter werden ausschließlich auf dem Developmentset bestimmt; anschließend werden fünf Finalisten für die Validierung eingefroren und der dort ausgewählte Sieger genau einmal auf dem zurückgehaltenen Testsplit evaluiert. Die ausgewählte Konfiguration aus multilingualem E5 und Top-1 mit Schwellenwert ($\theta=0.83959$) erreicht einen mittleren fragenweisen F1-Wert für die Chunkauswahl von $0.647$ auf dem Entwicklungs-, $0.611$ auf dem Validierungs- und $0.491$ auf dem Testsplit. Die vollständige Prüfung der 74 nicht exakt korrekten Testfälle identifiziert 21 Abrufe ähnlicher, aber falscher Chunks, 18 schwellenwertbedingte Enthaltungen bei beantwortbaren Fragen, 13 falsch-positive Abrufe bei Zero-Gold-Fragen, 9 strukturelle Top-1-Limitationen, 9 Fälle valider alternativer Chunks und 4 Probleme der Frageformulierung bzw. Lexik. Eine im Entwicklungs- und Validierungssplit ungesehene Evidenzgruppe mit acht Testfragen führt zu acht Fällen mit F1 $=0$ und verdeutlicht damit die strengere Generalisierung auf Evidenzgruppen, die durch die gegen Evidenz-Leckage abgesicherte Datenaufteilung entsteht.
\end{abstract}


\section{Einleitung}

\begin{leitfragenbox}
\begin{itemize}
    \item \textbf{Problem und Geltungsbereich:} Warum ist die pr\"azise Identifikation von Belegstellen in einer strukturierten Pr\"ufungsordnung ein eigenständiges Problem des Informationsretrievals? Welche Rolle spielen rechtlich nahe, aber funktional unterschiedliche Passagen, Tabellen und Studienverlaufsinformationen?
    \item \textbf{Retrieval als vorgelagerte Komponente:} Warum wird in dieser Arbeit isoliert \textit{Belegstellen-Retrieval} bzw. Vorhersage von Evidenzsets untersucht, statt generierte Antworten oder End-to-End-QA zu evaluieren? Wie ordnet sich Retrieval als vorgeschaltete Komponente wissensintensiver QA/RAG-Systeme ein, ohne daraus eine Garantie gegen Halluzinationen abzuleiten?
    \item \textbf{Enthaltung:} Warum muss das System f\"ur plausibel klingende, aber aus der Quelle nicht beantwortbare Fragen explizit $[\,]$ vorhersagen k\"onnen?
    \item \textbf{Forschungsfragen:}
    \begin{itemize}
        \item \textbf{RQ1 -- Repräsentation:} Wie unterscheiden sich spärliche lexikalische und dichte semantische Repräsentationen hinsichtlich ihrer Fähigkeit, antwortstützende Chunks aus einer deutschen Prüfungsordnung abzurufen?
        \item \textbf{RQ2 -- Selektion und Enthaltung:} Wie beeinflussen unterschiedliche Strategien zur Chunk-Selektion und Enthaltung die Retrieval-Leistung und das Enthaltungsverhalten?
        \item \textbf{RQ3 -- Generalisierung und Fehlermuster:} Wie gut generalisiert die ausgewählte Konfiguration auf den zurückgehaltenen Testsplit, und welche beobachteten Fehlermuster sowie Eigenschaften der Datenaufteilung und Evaluation stehen mit den verbleibenden Fehlern in Zusammenhang?
    \end{itemize}
    \item \textbf{Beiträge vorsichtig einordnen:} Empirischer Vergleich der drei Repr\"asentationen; separater Vergleich der Selektions-/Enthaltungsregeln; projekt-spezifischer kuratierter Retrieval-Datensatz; gegen Evidenzleckage abgesichertes Evaluationsdesign; Fehleranalyse des eingefrorenen Systems. Keine Behauptung einer erstmaligen Methode, keine SOTA-Behauptung und keine allgemeine Überlegenheitsbehauptung.
\end{itemize}
\end{leitfragenbox}

\begin{quellenbox}
\begin{itemize}
    \item \textbf{\citet{lewis2020rag}:} Prim\"arquelle daf\"ur, Retrieval als explizite nicht-parametrische Wissenskomponente in wissensintensiven NLP-/QA-Systemen zu motivieren. Nicht als Beleg verwenden, dass das hier untersuchte Retrieval Halluzinationen verhindert.
    \item \textbf{\citet{karpukhin-etal-2020-dense}:} Grundlegende Arbeit zum dichten Passage-Retrieval f\"ur offene domänenübergreifende QA; geeignet zur Einordnung von Bi-Encoder-Retrieval.
    \item \textbf{\citet{rajpurkar-etal-2018-know}:} Analogie f\"ur die wissenschaftliche Bedeutung unbeantwortbarer Fragen und expliziter Entscheidungen ohne Antwort.
\end{itemize}
\end{quellenbox}

\begin{faktenbox}
\begin{itemize}
    \item Gegenstand: Pr\"ufungsordnung 2026 f\"ur den BA Computerlinguistik im untersuchten Dokumentkorpus.
    \item Vorhersageziel: Menge von Retrieval-Chunk-IDs; keine generierte Antwortqualit\"at.
    \item Prim\"are Metrik: mittlerer fragenweiser F1 der Belegstellenauswahl.
    \item Offizieller eingefrorener Test-Endpunkt: F1 $=0.4907$ bei 144/144 eingereichten und ausgewerteten Vorhersagedatensätzen.
\end{itemize}
\end{faktenbox}

\textit{[Fließtextplatzhalter für Abschnitt 1]}



\section{Motivation und verwandte Arbeiten}

\begin{leitfragenbox}
\begin{itemize}
    \item \textbf{Sparsames vs. dichtes Retrieval:} Welche Grundlagen des Vektorraummodells und von TF-IDF sind f\"ur den Vergleich mit neuronalen Satz-/Passage-Embeddings relevant?
    \item \textbf{Allgemeine semantische Ähnlichkeit vs. retrieval-orientierte Kodierung:} Welche methodische Differenz besteht zwischen einem allgemeinen multilingualen Sentence-BERT-Modell-Checkpoint und E5 als retrieval-orientierter Text-Embedding-Familie? Nur als Motivation formulieren, nicht als kausale Erkl\"arung des Projektergebnisses.
    \item \textbf{Zero-Shot-Robustheit:} Was zeigt BEIR tats\"achlich? Insbesondere: BM25 bleibt auf heterogenen Zero-Shot-Retrieval-Aufgaben eine starke lexikalische Vergleichsverfahren. Daraus nicht ableiten, BEIR beweise die projektspezifische St\"arke von TF-IDF aufgrund administrativer Fachterminologie.
    \item \textbf{Selektive Vorhersage / Enthaltung:} Wie motivieren Ablehnungsoption und QA mit unbeantwortbaren Fragen das bewusste Zur\"uckweisen unsicherer bzw. nicht belegbarer Anfragen?
    \item \textbf{Kontext juristischer und regulatorischer Retrieval-Aufgaben:} Welche Arbeiten zeigen, dass juristische Retrieval-Aufgaben und strukturell \"ahnliche Rechtsdokumente besondere Evidenz-Disambiguierung erfordern?
    \item \textbf{Wissenschaftliche Einordnung:} Die Arbeit als kontrollierte Fallstudie in einem kleinen deutschen regulatorischen Korpus positionieren. Keine unbelegte Behauptung, erstmals Pr\"ufungsordnungs-Retrieval zu untersuchen.
\end{itemize}
\end{leitfragenbox}

\begin{quellenbox}
\begin{itemize}
    \item \textbf{Sparsames Retrieval:} \citet{sparckjones1972statistical} f\"ur Termspezifität/IDF und \citet{salton1975vector} f\"ur das klassische Vektorraummodell.
    \item \textbf{Satz-Embeddings:} \citet{reimers-gurevych-2019-sentence} f\"ur Sentence-BERT sowie \citet{reimers-gurevych-2020-making} f\"ur multilinguale Destillation multilingualer Satz-Embeddings.
    \item \textbf{Retrieval-Embeddings:} \citet{wang2022text} und \citet{wang2024multilingual} f\"ur die E5-Familie.
    \item \textbf{Benchmarks:} \citet{thakur2021beir} f\"ur heterogene Zero-Shot-IR-Evaluation und die Robustheit von BM25; \citet{muennighoff-etal-2023-mteb} f\"ur aufgabenübergreifende Embedding-Evaluation.
    \item \textbf{Enthaltung:} \citet{pmlr-v97-geifman19a} f\"ur selektive Vorhersage / Ablehnungsoption; \citet{rajpurkar-etal-2018-know} f\"ur QA mit unbeantwortbaren Fragen.
    \item \textbf{Deutsche Rechtsdom\"ane:} \citet{wrzalik-krechel-2021-gerdalir} und \citet{buttner-habernal-2024-answering} zur Einordnung deutscher juristischer IR-/QA-Aufgaben.
    \item \textbf{Aktueller Kontext für juristisches Retrieval:} \citet{reuter-etal-2025-towards} zeigt aktuelle Retrieval-Probleme in gro\ss en, strukturell \"ahnlichen juristischen Dokumentbest\"anden. Die Aufgabe ist nicht identisch mit dem hier untersuchten Einzelkorpus; die Quelle dient nur der Dom\"anenpositionierung.
\end{itemize}
\end{quellenbox}

\begin{faktenbox}
\begin{itemize}
    \item Verwandte Arbeiten kompakt halten; keine breite Historie des juristischen NLP.
    \item TF-IDF-Erkl\"arung im Ergebnisteil nur als \textit{plausible Erklärung}: hohe lexikalische Spezifit\"at wurde nicht als Kausalmechanismus experimentell isoliert.
    \item Sentence-BERT nur als im vorliegenden Versuchsaufbau schw\"achere Repr\"asentation beschreiben; keine Generalisierung auf juristisches NLP insgesamt.
\end{itemize}
\end{faktenbox}

\textit{[Fließtextplatzhalter für Abschnitt 2]}


\section{Methodik}

\begin{leitfragenbox}
\begin{itemize}
    \item \textbf{Textrepräsentationen:}
    \begin{itemize}
        \item \textit{TF-IDF:} Tats\"achlich ausgef\"uhrte Konfiguration: \texttt{TfidfVectorizer(lowercase=True, ngram\_range=(1,3))}. Das Vokabular wird auf den Chunk-Texten gefittet; die Fragen werden anschlie\ss end in denselben Raum transformiert. Keine sublineare TF-Skalierung behaupten, da \texttt{sublinear\_tf} nicht gesetzt wurde.
        \item \textit{Sentence-BERT:} Tats\"achlich verwendeter Modell-Checkpoint: \nolinkurl{sentence-transformers/paraphrase-multilingual-MiniLM-L12-v2} (384-dimensionale Satz-Embeddings). Kein selbst implementiertes Mittelwert-Pooling behaupten; die Repr\"asentation wird \"uber \texttt{SentenceTransformer} erzeugt.
        \item \textit{Multilingual E5:} Modell-Checkpoint \texttt{intfloat/multilingual-e5-base} (768 Dimensionen). Fragen werden mit \texttt{query:} und Chunks mit \texttt{passage:} pr\"afixiert.
        \item Keine der beiden dichten Repr\"asentationen wurde auf dem Projektkorpus feinabgestimmt.
    \end{itemize}
    \item \textbf{Ähnlichkeitsberechnung:} F\"ur jede Repr\"asentation wird Kosinusähnlichkeit zwischen allen 720 Fragen und 201 Chunks berechnet. Die resultierende Matrix hat $720\times201$ Eintr\"age. Bei den neuronalen Encodern wurde \texttt{normalize\_embeddings=False} verwendet; daher nicht behaupten, die gespeicherten dichten Vektoren seien vorab L2-normalisiert. Die Kosinusähnlichkeit selbst normiert durch die Vektornormen.
    \item \textbf{Selektorfamilien:}
    \begin{enumerate}
        \item \textit{Top-$k$:}
        \[
        S_{\mathrm{top}\text{-}k}(q)=\operatorname{argtop}_k\{s(q,c_i)\}.
        \]
        Gibt bei ausreichend großer Sammlung immer $k$ Chunks zur\"uck und kann allein nicht abstainieren.
        \item \textit{Absoluter Schwellenwert:}
        \[
        S_{\theta}(q)=\{c\in\mathcal{C}\mid s(q,c)\geq\theta\}.
        \]
        Kann $[\,]$ oder mehrere Chunks zur\"uckgeben.
        \item \textit{Top-$k$ + Schwellenwert:} Zunächst die Chunks, die den Schwellenwert überschreiten, nach Ähnlichkeitswert ordnen, dann h\"ochstens die ersten $k$ zur\"uckgeben. Damit werden absolute Filterung, Enthaltung und Begrenzung der Ausgabemenge kombiniert.
        \item \textit{Relative Marge:} Mit $s_k$ als Ähnlichkeitswert auf Rang $k$ und $s_{k+1}$ als folgendem Ähnlichkeitswert gilt
        \[
        m_k(q)=\frac{s_k-s_{k+1}}{\max(|s_k|,10^{-12})}.
        \]
        Top-$k$ wird nur ausgegeben, wenn $m_k(q)\geq\gamma$; andernfalls wird $[\,]$ vorhergesagt. F\"ur $k=1$ misst dies die relative Trennung zwischen Rang 1 und Rang 2.
    \end{enumerate}
    \item \textbf{Suchbereiche für Schwellenwerte:} Erkl\"aren, dass repr\"asentationsspezifische Regionen auf dem Entwicklungssplit aus robusten Quantilen von schlechtesten Gold-Scores, besten Nicht-Gold-Scores und höchsten Zero-Gold-Scores abgeleitet wurden; die eigentliche Parametersuche wurde ausschlie\ss lich auf dem Entwicklungssplit ausgewertet.
\end{itemize}
\end{leitfragenbox}

\begin{quellenbox}
\begin{itemize}
    \item \textbf{\citet{salton1975vector}:} Vektorraumdarstellung und Cosine-\"Ahnlichkeit als klassische IR-Grundlage.
    \item \textbf{\citet{reimers-gurevych-2019-sentence}, \citet{reimers-gurevych-2020-making}:} Sentence-BERT und multilinguale Satz-Embeddings.
    \item \textbf{\citet{wang2022text}, \citet{wang2024multilingual}:} Retrieval-orientiertes E5-Training und multilingualer E5-Kontext.
    \item \textbf{\citet{pmlr-v97-geifman19a}:} Konzeptionelle Grundlage f\"ur selektive Vorhersage; die konkrete Ähnlichkeits-Schwellenwertregel bleibt eine Designentscheidung dieses Experiments.
\end{itemize}
\end{quellenbox}

\begin{faktenbox}
\begin{itemize}
    \item Eingefrorener Vergleich im Paper: TF-IDF, Sentence-BERT, E5; Kosinusähnlichkeit; Top-$k$, Schwellenwert, Top-$k$+Schwellenwert, Relative Marge.
    \item Nicht in die Argumentation aufnehmen: im Code zus\"atzlich vorhandene Skalarprodukt-/euklidische Funktionen oder \texttt{softmax\_cumulative}, da sie nicht zu den eingefrorenen vier Selektor-Familien des Papers geh\"oren.
    \item Dokumentierte anfängliche Suchbereiche für Schwellenwerte: E5 $0.8331$--$0.8865$; TF-IDF $0.0591$--$0.5162$; Sentence-BERT $0.4740$--$0.7595$.
\end{itemize}
\end{faktenbox}

\textit{[Fließtextplatzhalter für Abschnitt 3]}


\section{Experimenteller Aufbau}

\begin{leitfragenbox}
\begin{itemize}
    \item \textbf{Korpus und Chunking:}
    \begin{itemize}
        \item Wie wurden die 201 aktiven Retrieval-Chunks aus allgemeiner Prüfungsordnung, fachspezifischem Anhang, exemplarischem Studienverlaufsplan und Vorspann gebildet?
        \item Welche Kontextüberschriften sorgen daf\"ur, dass kurze Tabellen-/Anhangswerte au\ss erhalb ihrer Originalseite interpretierbar bleiben?
        \item Welche finalen Qualit\"atschecks bestanden die Chunks (eindeutige IDs, keine leeren Chunks, vollst\"andige Quellreihenfolge)?
    \end{itemize}
    \item \textbf{Frage- und Goldstandard-Konstruktion:}
    \begin{itemize}
        \item 720 Fragen insgesamt: 600 beantwortbare und 120 Zero-Gold-Fragen.
        \item Gold-Ziel ist \texttt{all\_required\_chunk\_ids}: ein \textit{minimal vorgesehenes vollständiges Evidenzset}, nicht notwendigerweise die vollständige Menge aller semantisch gültigen Alternativbelege.
        \item \textbf{TODO -- Provenienz einfrieren:} Vor finaler Prosa exakt dokumentieren, welche Schritte bei Chunking, Fragenerstellung, Zuordnung und Pr\"ufung menschlich, skriptgest\"utzt bzw. ggf. sprachmodellassistiert waren. Keine unbelegte Darstellung als vollständig manuell oder vollständig LLM-generiert verwenden.
    \end{itemize}
    \item \textbf{Gegen Evidenzleckage abgesicherte Datenaufteilung:}
    \begin{itemize}
        \item Entwicklung/Validierung/Test: 432/144/144 Fragen.
        \item Gruppierungsregel \texttt{same\_non\_empty\_gold\_chunk\_set}: Fragen mit identischem nicht-leerem Gold-Evidenzset bleiben im selben Datenanteil.
        \item Integrität der Datenaufteilung: 315 Gruppen; gr\"o\ss te Gruppe 13 Fragen; keine nicht-leere Gold-Set-Gruppe in mehreren Datenanteilen vorkommt; Zero-Gold-Fragen bilden jeweils Einzelgruppen.
        \item Abwägung: weniger Evidenz-Leakage, dafür strengere Generalisierung auf vollständig ungesehene Evidenzgruppen.
    \end{itemize}
    \item \textbf{Selektionsprotokoll:}
    \begin{itemize}
        \item Selektor-/Hyperparameterentwicklung nur auf dem Entwicklungssplit.
        \item Ranking: mittlerer fragenweiser F1; dokumentierte Regeln bei Gleichstand: exakte Übereinstimmung, Präzision, dann weniger ausgew\"ahlte Chunks.
        \item F\"unf Finalisten vor Validierung einfrieren; keine Nachjustierungsschleife nach der Validierung; einen Sieger auf dem Validierungssplit einfrieren; genau diesen einmal auf dem Testsplit evaluieren.
    \end{itemize}
    \item \textbf{Evaluation:}
    \begin{itemize}
        \item $\text{Gold}=[\,]$ und $\text{Pred}=[\,]$ ergibt Präzision=Recall=F1=1 und exakte Übereinstimmung.
        \item $\text{Gold}=[\,]$, $\text{Pred}\neq[\,]$ sowie $\text{Gold}\neq[\,]$, $\text{Pred}=[\,]$ ergeben pro Frage F1=0.
        \item Ein komplett fehlender Vorhersagedatensatz ist semantisch anders als ein explizites $[\,]$: fehlende Datensätze werden aus den Mittelwerten ausgeschlossen und als unbeantwortet gez\"ahlt. Die eingefrorene Testauswertung weist 144/144 eingereicht/ausgewertet und 0 unbeantwortet aus.
        \item Sekundäre Metriken: mittlere fragenweise Präzision/Recall, exakte Übereinstimmung, Micro-F1, Zero-Gold-Enthaltungsrate, Rate leerer Selektionen, durchschnittliche Anzahl ausgewählter Chunks.
    \end{itemize}
    \item \textbf{Vergleichsverfahren:}
    \begin{itemize}
        \item \textit{Zufälliges Top-1}: gleichverteilt zufälliger Chunk; 100 feste Zufalls-Seeds 0--99; Mittelwert $\pm$ Standardabweichung berichten.
        \item \textit{Häufigste Entwicklungsantwort}: h\"aufigstes nicht-leeres exaktes Gold-Set auf dem Entwicklungssplit einmal einfrieren und f\"ur Validierung/Test konstant vorhersagen.
        \item \textit{Immer keine Auswahl}: f\"ur jede Frage $[\,]$.
        \item Autoritative Vergleichsverfahren auf dem Testsplit sind jetzt eingefroren: Zufälliges Top-1 (100 Zufalls-Seeds) erreicht mittleren Q-F1 $0.0047\pm0.0049$; Häufigste Entwicklungsantwort erreicht Q-F1 $0$; Immer keine Auswahl erreicht Q-F1 $0.1528$ und entspricht damit genau dem Anteil der 22 Zero-Gold-Fragen im 144-Fragen-Testsplit. Die vollst\"andigen Werte stehen in der Ergebnis-Faktenbox.
    \end{itemize}
\end{itemize}
\end{leitfragenbox}

\begin{quellenbox}
\begin{itemize}
    \item \textbf{\citet{rajpurkar-etal-2018-know}:} Literaturkontext f\"ur bewusst unbeantwortbare Fragen; die konkrete Zero-Gold-Konstruktion ist projektspezifisch.
    \item \textbf{\citet{muennighoff-etal-2023-mteb}:} Kontext f\"ur standardisierte Embedding-/Retrieval-Evaluation; die konkrete mittlere Set-F1-Definition des Projekts muss dennoch selbst explizit definiert werden.
\end{itemize}
\end{quellenbox}

\begin{faktenbox}
\begin{center}
\small
\textbf{Verifizierte Verteilung der Gold-Sets.}\\[1ex]
\resizebox{\linewidth}{!}{%
\begin{tabular}{lrrrr}
\toprule
\textbf{Split} & \textbf{Gesamt} & \textbf{Zero-Gold} & \textbf{1-Gold} & \textbf{2-Gold} \\
\midrule
Entwicklung & 432 & 74 & 340 & 18 \\
Validierung  & 144 & 24 & 116 & 4 \\
Test        & 144 & 22 & 112 & 10 \\
\midrule
\textbf{Gesamt} & \textbf{720} & \textbf{120} & \textbf{568} & \textbf{32} \\
\bottomrule
\end{tabular}}
\end{center}

\vspace{1ex}
\begin{itemize}
    \item 201 aktive Retrieval-Chunks; 10,194 W\"orter; mittlere Chunkl\"ange 50.72 W\"orter; Median 41; Spannweite 8--204 W\"orter.
    \item Verteilung der Quellschichten aus dem Qualitätsaudit: 133 Allgemeine Prüfungsordnung, 36 Exemplarischer Studienverlaufsplan, 30 Fachspezifischer Anhang, 2 Vorspann.
    \item Die Chunk-Sammlung ist als \texttt{ready\_with\_minor\_issues} auditiert; dokumentierte kleinere Quell-/Formatrisiken nicht als fehlerhafte Vorverarbeitung darstellen.
    \item \textbf{Eingefrorene Reproduzierbarkeitsumgebung:} Projektprotokolle verwenden Python 3.11; der gepinnte PyTorch-Paket ist ebenfalls f\"ur CPython 3.11 (\texttt{cp311}). Materiell relevante Bibliotheksversionen aus \texttt{requirements.txt}: \texttt{scikit-learn==1.9.0}, \texttt{sentence-transformers==3.4.1}, \texttt{transformers==4.46.3}, \texttt{torch==2.1.2} (CPU-Build), \texttt{numpy==1.26.4}, \texttt{scipy==1.17.1}, \texttt{pandas==3.0.5}, \texttt{tokenizers==0.20.3} und \texttt{huggingface-hub==0.36.2}. In der finalen Prosa gen\"ugen die f\"ur Reproduzierbarkeit relevanten Kernpakete; die komplette Abhängigkeitsliste bleibt im Projektartefakt \texttt{requirements.txt}.
\end{itemize}
\end{faktenbox}

\textit{[Fließtextplatzhalter für Abschnitt 4]}


\section{Ergebnisse und Evaluation}

\begin{leitfragenbox}
\begin{itemize}
    \item \textbf{RQ1 -- Vergleich der Repräsentationen:}
    \begin{itemize}
        \item E5 erreicht auf dem Entwicklungssplit die st\"arksten Werte unter den evaluierten Repr\"asentationen.
        \item TF-IDF bleibt in diesem Korpus konkurrenzf\"ahig. Zul\"assige Interpretation: Die hohe lexikalische Spezifit\"at des Korpus ist eine plausible Erkl\"arung; sie wurde nicht als Kausalmechanismus isoliert gemessen.
        \item Der evaluierte multilingual MiniLM Sentence-BERT-Modell-Checkpoint liegt im betrachteten Versuchsaufbau hinter E5 und TF-IDF. Keine Generalisierung wie ``Sentence-BERT funktioniert nicht f\"ur Rechtstexte''.
    \end{itemize}
    \item \textbf{RQ2 -- Selektor und Enthaltung:}
    \begin{itemize}
        \item Nicht nach einem universell ``\"uberlegenen'' Selektor fragen, sondern nach Abwägungen zwischen Belegstellen-Retrieval, Begrenzung der Ausgabemenge und Enthaltung.
        \item Reines Top-1 ist stark, kann aber nie abstainieren. Reiner Schwellenwert kann abstainieren, kann jedoch viele Chunks selektieren und dadurch Präzision verlieren.
        \item Top-1 + absoluter Schwellenwert erzeugt die st\"arkste E5-Konfiguration auf dem Entwicklungssplit und bleibt der Sieger auf dem Validierungssplit; Relative Marge ist konkurrenzf\"ahig, insbesondere f\"ur E5/TF-IDF, aber nicht der finale Sieger auf dem Validierungssplit.
        \item Wichtig: F\"ur TF-IDF ist Relative Marge auf dem Entwicklungssplit geringf\"ugig st\"arker als Top-1+Schwellenwert; daher keine universelle Selektor-Rangordnung behaupten.
    \end{itemize}
    \item \textbf{Validierung und eingefrorene Modellauswahl:}
    \begin{itemize}
        \item E5 + Top-1 + Schwellenwert $0.83959$: F1 $=0.6111$, EM $=0.5972$, mittlere Präzision $=0.6181$, mittlerer Recall $=0.6076$, Zero-Gold-Enthaltungsrate $=0.4167$.
        \item Nur die f\"unf eingefrorenen Finalisten wurden auf dem Validierungssplit verglichen; Sentence-BERT war nicht unter diesen f\"unf Finalisten. Daher Validierung nicht als erneuten vollst\"andigen Drei-Repr\"asentationen-Vergleich formulieren.
    \end{itemize}
    \item \textbf{RQ3 -- Zurückgehaltener Testsplit:}
    \begin{itemize}
        \item Offizieller Test: F1 $=0.4907$, EM $=0.4861$, mittlere Präzision $=0.4931$, mittlerer Recall $=0.4896$, Micro-F1 $=0.5041$.
        \item Rate leerer Selektionen $=0.2083$; Zero-Gold-Enthaltung $=9/22=0.4091$. Das System zeigt sowohl falsche Enthaltungen bzw. zu geringe Retrieval-Ausgabe als auch falsch-positive Abrufe bei Zero-Gold.
        \item Rückgang von Validierung zu Test: ca. $0.1204$. Keine Aussage, dieser Rückgang sei ``innerhalb einer normalen Standardabweichung'', solange keine entsprechende Unsicherheitsanalyse durchgef\"uhrt wurde.
    \end{itemize}
    \item \textbf{Fehleranalyse:}
    \begin{enumerate}
        \item \textit{Strukturelle Limitation bei mehreren Belegstellen:} Bei echten 2-Gold-Fragen kann ein Top-1-System nie eine exakte Übereinstimmung des vollständigen Sets erzielen. Wenn genau einer der zwei Gold-Chunks korrekt geliefert wird, ist maximal $P=1$, $R=0.5$, $F1=2/3$. Im Test gibt es 10 2-Gold-Fragen.
        \item \textit{Enthaltungen bei beantwortbaren Fragen} vs. \textit{falsch-positive Abrufe bei Zero-Gold} getrennt quantifizieren.
        \item \textit{Fall einer ungesehenen Evidenzgruppe:} Gruppe G0040 umfasst acht Testfragen Q0153--Q0160 mit demselben Gold-Chunk \texttt{PO25CL-GEN-C03-P12-A02}; dieses nicht-leere Gold-Set kommt in Entwicklung und Validierung nicht vor. Alle acht Fragen erzielen auf dem Testsplit F1 $=0$: Q0153, Q0154, Q0155, Q0156, Q0158 und Q0159 liefern jeweils einen falschen nicht-leeren Chunk; Q0157 und Q0160 abstainieren mit $[\,]$. Der Cluster ist daher ein verifiziertes Beispiel f\"ur Generalisierungsfehler auf Evidenzgruppen, nicht f\"ur einen fehlerhaften Testsplit.
        \item \textit{Strenge des Goldstandards:} Die abgeschlossene Prüfung falscher Chunks klassifiziert 9 formal falsche Chunk-ID-Vorhersagen als \texttt{valid\_alternative\_evidence}. Diese F\"alle zeigen, dass das minimal vorgesehene Gold-Set nicht jede semantisch vollständig ausreichende Alternativstelle enthält; der offizielle offizielle Wert bleibt unver\"andert.
        \item \textit{Vollständige Fehlerprüfung:} Die 31 Fälle mit falscher nicht-leerer Vorhersage und 43 strukturellen bzw. enthaltungsbezogenen Fälle decken zusammen alle 74 nicht exakt korrekten Testfragen ab. Finale Kategorien: 21 ähnliche, aber falsche Chunks, 18 schwellenwertbedingte Enthaltungsfehler, 13 falsch-positive Abrufe bei unbeantwortbaren Fragen, 9 strukturelle Top-1-Limitationen, 9 Fälle valider alternativer Belegstellen und 4 Probleme der Frageformulierung bzw. Lexik.
        \item \textit{Post-hoc-Diagnose (optional):} Da nun sowohl die 9 Fälle mit valider alternativer Evidenz als auch alle acht Fälle mit F1 = 0 der G0040-Gruppe verifiziert sind, kann eine klar als Sensitivit\"atsrechnung bezeichnete Variante berechnet werden:
        \[
        \frac{70.6667 + 9}{144-8} \approx 0.5858.
        \]
        Dieser Wert ist \emph{kein} korrigierter Testwert und darf nicht in die Hauptresultattabelle eingehen; er zeigt lediglich, wie stark diese beiden identifizierten Daten-/Evaluationsph\"anomene den rohen Leistungsabstand beeinflussen k\"onnen.
    \end{enumerate}
    \item \textbf{Vergleich mit Vergleichsverfahren:} Die eingefrorenen Vergleichsverfahren auf dem Testsplit sind verifiziert und werden in der Ergebnis-Faktenbox berichtet.
\end{itemize}
\end{leitfragenbox}

\begin{quellenbox}
\begin{itemize}
    \item \textbf{\citet{thakur2021beir}:} Nur zur allgemeinen Einordnung starker lexikalischer Zero-Shot-Vergleichsverfahren; nicht als Beweis der projektspezifischen TF-IDF-Ursache.
    \item \textbf{\citet{reuter-etal-2025-towards}:} Aktueller Dom\"anenkontext daf\"ur, dass Retrieval in strukturell \"ahnlichen juristischen Best\"anden fehleranf\"allig sein kann; keine direkte Erkl\"arung einzelner Projektfehler.
    \item \textbf{Eigene eingefrorene Evaluation:} Alle Zahlen und Fehleranzahlen aus Projektartefakten berichten; Literatur darf diese Messergebnisse nicht ersetzen.
\end{itemize}
\end{quellenbox}

\begin{faktenbox}
\begin{center}
\small
\textbf{Entwicklung: beste Konfigurationen je Repr\"asentation und Selektor-Familie. Alle Präzisions-/Recall-Spalten enthalten mittlere fragenweise Metriken.}\\[1ex]
\resizebox{\linewidth}{!}{%
\begin{tabular}{llcrrrr}
\toprule
\textbf{Repräsentation} & \textbf{Selektor} & \textbf{Parameter} & \textbf{Q-F1} & \textbf{EM} & \textbf{Präzision} & \textbf{Recall} \\
\midrule
\textbf{E5} & \textbf{Top-1 + Schwelle} & $\theta=0.83959$ & \textbf{0.6474} & \textbf{0.6366} & \textbf{0.6528} & \textbf{0.6447} \\
E5 & Relative Marge & $\gamma=0.0025$ & 0.6312 & 0.6204 & 0.6366 & 0.6285 \\
E5 & Top-1 & $k=1$ & 0.6165 & 0.6042 & 0.6227 & 0.6134 \\
E5 & Schwellenwert & $\theta=0.866475$ & 0.4941 & 0.3657 & 0.4578 & 0.6285 \\
\midrule
TF-IDF & Relative Marge & $\gamma=0.235$ & 0.5702 & 0.5579 & 0.5764 & 0.5671 \\
TF-IDF & Top-1 + Schwelle & $\theta=0.0999$ & 0.5671 & 0.5532 & 0.5741 & 0.5637 \\
TF-IDF & Top-1 & $k=1$ & 0.5509 & 0.5370 & 0.5579 & 0.5475 \\
TF-IDF & Schwellenwert & $\theta=0.127665$ & 0.4866 & 0.3657 & 0.4596 & 0.5706 \\
\midrule
Sentence-BERT & Top-1 + Schwelle & $\theta=0.6751$ & 0.3989 & 0.3958 & 0.4005 & 0.3981 \\
Sentence-BERT & Relative Marge & $\gamma=0.005$ & 0.3866 & 0.3819 & 0.3889 & 0.3854 \\
Sentence-BERT & Top-1 & $k=1$ & 0.3742 & 0.3681 & 0.3773 & 0.3727 \\
Sentence-BERT & Schwellenwert & $\theta=0.7095375$ & 0.3583 & 0.2801 & 0.3343 & 0.4421 \\
\bottomrule
\end{tabular}}
\end{center}

\vspace{1.5ex}
\begin{center}
\small
\textbf{Eingefrorener Sieger: Validierung und zurückgehaltener Test.}\\[1ex]
\begin{tabular}{lrrrrrr}
\toprule
\textbf{Split} & \textbf{Q-F1} & \textbf{EM} & \textbf{Präz.} & \textbf{Rec.} & \textbf{Micro-F1} & \textbf{ZG-Enth.} \\
\midrule
Validierung & 0.6111 & 0.5972 & 0.6181 & 0.6076 & 0.6475 & 0.4167 \\
Test       & 0.4907 & 0.4861 & 0.4931 & 0.4896 & 0.5041 & 0.4091 \\
\bottomrule
\end{tabular}
\end{center}

\vspace{1ex}
\textbf{Eingefrorene Details des Testsystems:} 144 eingereicht / 144 ausgewertet / 0 unbeantwortet; durchschnittliche Anzahl ausgewählter Chunks $=0.7917$; Rate leerer Selektionen $=0.2083$; TP=62, FP=52, FN=70.

\vspace{1.5ex}
\begin{center}
\small
\textbf{Zurückgehaltenes Testsystem im Vergleich zu den eingefrorenen Vergleichsverfahren. Für Zufälliges Top-1 werden Mittelwert $\pm$ Standardabweichung über 100 feste Zufalls-Seeds berichtet.}\\[1ex]
\begin{tabular}{lrrrrr}
\toprule
\textbf{System} & \textbf{Q-F1} & \textbf{Präz.} & \textbf{Rec.} & \textbf{EM} & \textbf{ZG-Enth.} \\
\midrule
Zufälliges Top-1 & $0.0047\pm0.0049$ & $0.0050\pm0.0050$ & $0.0045\pm0.0048$ & $0.0041\pm0.0048$ & 0.0000 \\
Häufigste Entwicklungsantwort & 0.0000 & 0.0000 & 0.0000 & 0.0000 & 0.0000 \\
Immer keine Auswahl & 0.1528 & 0.1528 & 0.1528 & 0.1528 & 1.0000 \\
\textbf{E5 Top-1 + Schwellenwert} & \textbf{0.4907} & \textbf{0.4931} & \textbf{0.4896} & \textbf{0.4861} & \textbf{0.4091} \\
\bottomrule
\end{tabular}
\end{center}

\vspace{1.5ex}
\begin{center}
\small
\textbf{Eingefrorene Kategorisierung aller 74 nicht exakt korrekten Testfälle.}\\[1ex]
\begin{tabular}{lrr}
\toprule
\textbf{Fehlerkategorie} & \textbf{Anzahl} & \textbf{Anteil an nicht exakt korrekten Fällen} \\
\midrule
Ähnlicher, aber falscher Chunk & 21 & 28.4\% \\
Schwellenwertbedingte Fehl-Enthaltung & 18 & 24.3\% \\
Falsch-positiver Abruf bei unbeantwortbarer Frage & 13 & 17.6\% \\
Strukturelle Top-1-Limitation & 9 & 12.2\% \\
Valide alternative Belegstelle & 9 & 12.2\% \\
Frageformulierung / lexikalisches Problem & 4 & 5.4\% \\
\midrule
\textbf{Gesamt} & \textbf{74} & \textbf{100.0\%} \\
\bottomrule
\end{tabular}
\end{center}
\end{faktenbox}

\textit{[Fließtextplatzhalter für Abschnitt 5]}


\section{Fazit}

\begin{leitfragenbox}
\begin{itemize}
    \item \textbf{RQ1 direkt beantworten:} E5 erzielt im evaluierten Datensatz die st\"arksten Ergebnisse auf Entwicklungs- und Validierungssplit; TF-IDF bleibt konkurrenzf\"ahig; der evaluierte Sentence-BERT-Modell-Checkpoint ist schw\"acher. Keine Generalisierung \"uber diesen Datensatz hinaus.
    \item \textbf{RQ2 direkt beantworten:} Repräsentation und Selektor tragen beide wesentlich zum gemessenen Verhalten bei. Der finale E5 Top-1+Schwellenwert-Selektor verbindet Ranking, Ausgabegrenze und Enthaltung und ist die auf dem Validierungssplit ausgewählte Konfiguration; daraus keine universelle Dominanz von Top-1+Schwellenwert f\"ur jede Repr\"asentation ableiten.
    \item \textbf{RQ3 direkt beantworten:} Die eingefrorene Konfiguration generalisiert schw\"acher auf dem Testsplit als auf dem Validierungssplit. Die vollst\"andige Auditierung der 74 nicht exakt korrekten Fälle zeigt mehrere gleichzeitig wirksame Fehlermuster: Abruf falscher Chunks, schwellenwertbedingte Enthaltung, falsch-positiver Abruf bei Zero-Gold, die strukturelle Top-1-Grenze, valide alternative Belegstelle sowie Frageformulierungsprobleme. Die acht Fälle mit F1 = 0 der ungesehenen Evidenzgruppe G0040 sind ein besonders deutliches Beispiel f\"ur die strengere Generalisierung auf Evidenzgruppen.
    \item \textbf{Praktische Implikation:} Belegstellen-Retrieval kann als klar evaluierbare vorgelagerte Komponente f\"ur sp\"atere QA/RAG-Systeme dienen; diese Arbeit selbst evaluiert jedoch keine generierten Antworten.
\end{itemize}
\end{leitfragenbox}

\textit{[Fließtextplatzhalter für Abschnitt 6]}


\section{Limitationen und Ausblick}

\begin{leitfragenbox}
\begin{itemize}
    \item \textbf{Korpus- und Dom\"aneneinschr\"ankung:} Ein einzelnes regulatorisches Dokument / ein Studiengang; Generalisierbarkeit auf andere Pr\"ufungsordnungen, Hochschulen oder Rechtsdokumente ist nicht gezeigt.
    \item \textbf{Datensatzerstellung:} Projekt-spezifisch konstruierte Fragen und minimale Gold-Evidenzsets; Provenienz und Prüfprozess m\"ussen transparent berichtet werden. Keine Inter-Annotator-Reliabilität behaupten, falls keine unabh\"angige Doppelannotation vorliegt.
    \item \textbf{Harte Top-1-Begrenzung:} 10 Testfragen ben\"otigen zwei Gold-Chunks; der finale Selektor kann h\"ochstens einen Chunk liefern und ist damit f\"ur exaktes Retrieval des vollständigen Sets strukturell begrenzt.
    \item \textbf{Kalibrierung der Enthaltung:} Auf Test werden nur 9 von 22 Zero-Gold-Fragen korrekt abstainiert; zugleich entstehen Enthaltungen bei beantwortbaren Fragen. Das System ist daher weder einfach nur zu konservativ noch nur zu freizügig.
    \item \textbf{Geltungsbereich des Goldstandards:} Das Gold-Mapping bezeichnet ein minimal vorgesehenes Evidenzset. Die finale Prüfung falscher Chunks identifiziert 9 F\"alle, in denen eine formal falsche Chunk-ID dennoch voll ausreichende alternative Belegstelle liefert. Das ist eine Limitation der Exhaustivit\"at des Goldstandards, nicht ein Grund, den offiziellen Testwert nachtr\"aglich zu ersetzen.
    \item \textbf{Nicht evaluierte Methoden:} BM25, hybride Sparse-Dense-Verfahren, Cross-Encoder-Reranking, adaptive Auswahl mehrerer Chunks und generative Lesemodelle wurden nicht im eingefrorenen Vergleich getestet.
    \item \textbf{Weiterführende Arbeiten:} adaptives Retrieval mehrerer Belegstellen, BM25/hybrides Retrieval, Reranking, inhaltstypbewusstes Retrieval, breitere regulatorische Korpora, Alias-/Kanonisierungsunterstützung sowie nachgelagerte RAG-Evaluation nach abgeschlossenem Retrieval.
\end{itemize}
\end{leitfragenbox}

\begin{quellenbox}
\begin{itemize}
    \item \textbf{\citet{lewis2020rag}:} Geeigneter Literaturanker f\"ur sp\"atere nachgelagerte RAG-Evaluation nach dem Retrieval-Schritt.
    \item \textbf{\citet{reuter-etal-2025-towards}:} Aktueller Kontext für juristisches Retrieval f\"ur robuste Retrieval-/Chunking-Fragestellungen; nicht als Beleg daf\"ur verwenden, dass die dort vorgeschlagene Methode im vorliegenden Projekt getestet wurde.
    \item \textbf{\citet{thakur2021beir}:} Benchmark-Kontext f\"ur weitere lexikalische/dichte Retrieval-Vergleiche; nicht als Prim\"arquelle f\"ur ein konkretes hybrides Verfahren ausgeben.
\end{itemize}
\end{quellenbox}

\begin{faktenbox}
\begin{itemize}
    \item \textbf{Offener Punkt vor dem Einfrieren des Manuskripts:} Die Provenienz der Datensatzerstellung als transparenten Absatz ergänzen.
\end{itemize}
\end{faktenbox}

\textit{[Fließtextplatzhalter für Abschnitt 7]}

\bibliography{references}

\end{document}
